\documentclass[10pt,a4paper]{article}
\usepackage[utf8]{inputenc}
\usepackage[T1]{fontenc}
\usepackage[margin=18mm]{geometry}
\usepackage{longtable}
\usepackage{hyperref}
\usepackage{enumitem}
\setlist[itemize]{noitemsep, topsep=2pt}
\setlength{\parindent}{0pt}
\begin{document}

\section*{Non-Verbal \& Multimodal Discourse Analysis Worksheet}

\textit{Qualitative analysis of pauses, trailing sentences, prosody/intonation, vocalisations (coughs, sighs, scoffs), and kinesics (hand gestures, facial expressions, posture) in political news discourse}\par\smallskip

\bigskip\hrule\bigskip

\subsection*{Guidelines: Five non-verbal \& paralinguistic features to look for in every event}

\begin{longtable}{|p{0.23\textwidth}|p{0.23\textwidth}|p{0.23\textwidth}|p{0.23\textwidth}|}
\hline
Feature & What to look for & Illustrative example (from sample) & Questions to answer \\ \hline \endhead
\textbf{1. Pauses, Silences \& Abrupt Stops} & Micro-pauses \texttt{(.)}, timed silences \texttt{(1.5)}, pregnant/dramatic pauses before a loaded word, hesitant stumbles, abrupt self-interruptions or cut-offs (\texttt{---}), and forced silences (muting a panelist's mic or talking over them). & Anchor pauses \texttt{(1.5)} before saying \textit{"'loota'"} for dramatic weight; Addl. DCP cuts off mid-sentence (\textit{"wo toh--- abhi investigation chal raha hai"}) when asked if police struck first; panelist's mic faded during counter-question. & Where does the pause or stop occur? Is it deliberate/rhetorical (for suspense, emphasis, or gravitas) or defensive/evasive (cognitive load, self-repair, avoiding blame)? Who controls turn-taking and silence? \\ \hline
\textbf{2. Trailing Sentences \& Unfinished Utterances} & Clauses that trail off (\texttt{...}), drop in volume at the end, or leave the conclusion unspoken (aposiopesis) so the audience supplies the implication without the speaker taking explicit responsibility. & \textit{"Aur phir wahan jo log ghuse the... aap samajh hi sakte hain..."}; \textit{"500 salon ke sangharsh ke baad agar aisa..."} trailing off with a slow head shake. & What exact thought is left unfinished? What presupposition or shared bias is the viewer invited to fill in? Does trailing off act as a safe hedge to imply higher-level complicity or "anti-national" motive without stating it? \\ \hline
\textbf{3. Intonation, Pitch, Volume \& Prosody} & Pitch contour (rising \texttt{$\uparrow$} for incredulity/challenge; falling \texttt{$\downarrow$} for finality), volume shifts (shouting vs. hushed whisper), emphatic word stress, speech rate (rapid-fire urgency vs. slow solemnity), and tonal register (sarcastic, reverential, combative). & Hushed, reverential cadence on \textit{"Prabhu Shri Ram ke Mandir"} vs. sharp, high-pitched interrogation on \textit{"par Akhilesh ji Ram Mandir kyon$\uparrow$ nahin jate?!"}; satirical sing-song inflection on \textit{"priventiv arenjment"}. & Which exact words receive pitch peaks or volume stress? How does vocal tone turn neutral words into irony, mockery, or moral outrage? Does the speaker shift prosodic register depending on which actor is discussed? \\ \hline
\textbf{4. Coughs, Sighs \& Non-Lexical Vocalisations} & Audible sighs (\texttt{[sigh]}), scoffs/snorts (\texttt{[scoff]}), ironic chuckles (\texttt{[chuckle]}), throat-clearing/coughs (\texttt{[cough]}), sharp intakes of breath/gasps (\texttt{[gasp]}), and tongue-clicks (\texttt{[tsk]}). & Heavy \texttt{[sigh]} before reporting Trust oversight lapses (signalling reluctant sorrow rather than hostility); mocking \texttt{[scoff]} or \texttt{[chuckle]} when reading an opponent's denial; \texttt{[cough / throat-clear]} when pivoting away from an awkward fact. & What emotion or epistemic stance does the vocalisation perform (exasperation, moral disgust, disbelief, ridicule, discomfort)? Does a sigh or scoff pre-judge a speaker's claim before any verbal rebuttal is made? \\ \hline
\textbf{5. Hand Gestures, Facial Expressions \& Posture (Kinesics)} & Hand movements (air-quotes, finger-pointing/wagging, desk-thumping, open palms, counting on fingers, dismissive wave), facial cues (raised eyebrows, eye-rolls, smirk, furrowed brow), head nods/shakes, leaning into camera, and visual props (bandaged limbs, documents). & Anchor forms physical air-quotes on \textit{"'loota'"} or \textit{"'peaceful' protest"}; host points finger/pen at lens during an ultimatum; injured DCP holds bandaged hand up to camera; creator raises one eyebrow and leans in for a conspiratorial reveal. & How do gestures reinforce, contradict, or sharpen the spoken words? During split-screen listening, does the host's face (nodding vs. smirking/eye-rolling) cue the audience on whom to trust? How are bodily visuals used as evidence? \\ \hline
\end{longtable}

\subsection*{Quick Multimodal Transcription Notation Key (use inside transcript excerpts)}

\begin{longtable}{|p{0.15\textwidth}|p{0.15\textwidth}|p{0.15\textwidth}|p{0.15\textwidth}|p{0.15\textwidth}|p{0.15\textwidth}|}
\hline
Symbol & Meaning & Example in transcript & Symbol & Meaning & Example in transcript \\ \hline \endhead
\texttt{(.)} & Brief micro-pause (< 0.5 sec) & \texttt{woh (.) suparavijan lacking tha} & \texttt{$\uparrow$ / $\downarrow$} & Marked rising / falling pitch & \texttt{kyon$\uparrow$ FIR nahin karaee$\downarrow$} \\ \hline
\texttt{(1.5)} & Timed pause in seconds & \texttt{kisne (1.5) 'loota'?} & \texttt{WORD / \_word\_} & Loud volume / emphatic stress & \texttt{ye PAAP hai / \_bilkul\_ galat} \\ \hline
\texttt{---} & Abrupt stop, cut-off, or self-repair & \texttt{police ki taraf se--- abhi jaanch...} & \texttt{$^\circ$word$^\circ$} & Softened / whispered delivery & \texttt{$^\circ$ander ki baat ye hai$^\circ$} \\ \hline
\texttt{...} & Trailing off / unfinished sentence & \texttt{aap jaante hain ki kaun log...} & \texttt{[sigh] [cough] [scoff]} & Non-lexical vocalisations & \texttt{[sigh] 500 salon ke baad...} \\ \hline
\texttt{[overlap]} & Simultaneous / overlapping speech & \texttt{[Anjana: suniye---] [Guest: nahin!]} & \texttt{\{\{gesture: ...\}\}} & Hand gesture / facial cue / posture & \texttt{\{\{air-quotes on 'loota'\}\}} \\ \hline
\end{longtable}

\textbf{Methodological Reminder:} A non-verbal cue is not automatically manipulative or ideological---a cough, pause, or rapid cadence may reflect studio audio delay, live time pressure, or a speaker's habitual style. Always check (1) whether the cue clusters with verbal features, and (2) whether the speaker's delivery shifts noticeably across different actors or events (within-outlet baseline).\par\smallskip

\bigskip\hrule\bigskip

\subsection*{Event 1. Ram Mandir donation controversy (Non-Verbal \& Paralinguistic Layer)}

\begin{itemize}
  \item \textbf{Focus:} How do prosody, pauses, trailing sentences, vocalisations (sighs, coughs, scoffs), and body language frame the controversy---as solemn devotional sorrow, administrative scrutiny, or partisan combat?
  \item \textbf{Outlet:} Sarthak Goswami (\textit{Sunday Show}) | \textbf{Date:} 23 June 2021
  \item \textbf{Video title:} Ram Mandir Scam, Waqf Amendment Bill, Coaching Mafia, Floods \& Amritpal | \href{https://youtu.be/lb7gZnyZH0c}{https://youtu.be/lb7gZnyZH0c} | \textbf{Timestamp(s):} \texttt{[00:00]}--\texttt{[05:38]} (Full video: \texttt{19:33}; Ram Mandir segment: \texttt{[00:00]}--\texttt{[05:38]})
  \item \textbf{Analyst:} Baani Solanki (\texttt{2025114016}) | \textbf{Language:} Hinglish (Hindi in Roman script + English terms)
  \item \textbf{Format / Visual setup:} \texttt{[ ]} TV Panel Debate \texttt{[ ]} TV Bulletin \texttt{[x]} Solo YouTube Explainer \texttt{[ ]} Interview / Discussion
  \item \textbf{Primary speaker(s) observed:} Sarthak Goswami (solo studio presenter seated at desk, speaking directly to camera with jump-cut editing, meme inserts, and document overlays)
\end{itemize}

\subsubsection*{Event-specific non-verbal questions (Event 1)}

\noindent\textbf{1.} \textbf{Vocal \& gestural marking of core labels ("scam", "loota", "paap", "gadabdi", "avyavastha"): When the controversy is named, what non-verbal cues accompany the label? (Look for physical air-quotes, vocal scare-quotes, dramatic pauses before the word, hushed/reverential drops in volume, or shouted stress.)}\par\smallskip
\begin{itemize}
  \item \textbf{Notes:}
Sarthak Goswami avoids both the hushed, mournful devotional register (\texttt{"paap"}) and high-decibel shouting (\texttt{"scam!"}). Instead, when introducing the controversy at \texttt{[01:13]}, he uses a brief micro-pause \texttt{(.)} followed by raised eyebrows and open, upturned palms \texttt{\{\{gesture: open palms raised with arched eyebrows\}\}} on \textit{"Ayodhya Ram Mandir ke naam par (.) \_bada\_ ghotala$\uparrow$"} (\texttt{[01:13]}). Rather than endorsing \textit{"ghotala"} as his own verdict, his rising pitch \texttt{$\uparrow$} and wide-eyed expression perform \textbf{vocal scare-quotes} that mark the label as a dramatic headline circulating in public discourse. Similarly, when asking \textit{"kya sach mein Ram Mandir ke chande mein chori hui hai$\uparrow$"} (\texttt{[01:16]}), he leans slightly forward with an inquisitive head tilt, framing \textit{"chori"} as an open investigative question rather than an established moral sin.\par\smallskip
\end{itemize}

\noindent\textbf{2.} \textbf{Differential demeanour toward actors (Temple, Trust leadership, counting clerks, donors, opposition accusers): How do the speaker's voice pitch, facial expression, and posture shift when referring to Prabhu Shri Ram / devotees vs. top figures (Champat Rai, Nripendra Misra, CM Yogi) vs. arrested counting staff vs. opposition leaders (Akhilesh Yadav, Awadhesh Prasad)?}\par\smallskip
\begin{itemize}
  \item \textbf{Notes:}
Sarthak maintains a remarkably uniform, conversational baseline across all political and institutional actors, punctured equally by deadpan comedic irony:\par\smallskip
  \item \textbf{Donors / Devotees:} Referred to with warm, relatable facial animation and an inclusive hand gesture toward the camera (\texttt{"aapne aur humne milkar chanda diya tha"}, \texttt{[01:36]}), treating viewers as fellow citizens who contributed in good faith.
  \item \textbf{Opposition Accusers (Sanjay Singh, Pawan Pandey):} Introduced in a brisk, matter-of-fact narrative cadence (\texttt{[01:55]}), pointing to screen graphics displaying their press conference clips without scowling or mocking them.
  \item \textbf{Trust Leadership (Champat Rai):} When quoting Champat Rai's counter-attack comparing critics to \textit{"Mahatma Gandhi ke hatyare"} (\texttt{[04:10]}), Sarthak pauses \texttt{(1.0)}, deadpans to the camera with a raised eyebrow and an ironic half-smirk \texttt{\{\{gesture: deadpan stare into lens + single raised eyebrow\}\}}, letting the disproportionate hyperbole of the Trust's defence speak for itself.
  \item \textbf{Intermediaries (Sultan Ansari, Ravi Mohan Tiwari):} Listed with methodical finger-counting (\texttt{[02:10]}, \texttt{[02:27]}) and steady eye contact, treating them as transactional actors in a real-estate flip rather than demonised villains.
\end{itemize}

\noindent\textbf{3.} \textbf{Solemn devotional grief vs. high-decibel partisan spectacle: How do pauses, sighs, speech rate, and hand gestures construct the segment's emotional register? Where is delivery slow, hushed, and mournful (sighs, bowed head, folded hands), and where does it escalate into rapid-fire shouting, desk-thumping, or mocking irony?}\par\smallskip
\begin{itemize}
  \item \textbf{Notes:}
Sarthak explicitly rejects both solemn devotional grief and high-decibel partisan spectacle, operating instead in a \textbf{fast-paced, satirical infotainment register} (\textasciitilde{}155--165 words per minute). At \texttt{[00:34]}, he explicitly parodies TV news histrionics: \textit{"TV news kholte hi lagta hai (.) duniya \_khatam\_ hone wali hai"} (\texttt{[00:34]}), widening his eyes in mock panic and waving both hands before dropping back to a relaxed smile and level tone (\texttt{"isiliye shanti se samajhte hain"}, \texttt{[00:39]}). There are zero mournful bowed heads or folded hands, and zero desk-thumping outbursts. Emotional peaks are replaced by dry, punchy comedic pauses---such as a \texttt{(1.0)} beat before delivering the satirical line \textit{"zameen ki keemat (.) \_rocket\_ ki speed se badh gayi$\uparrow$"} (\texttt{[02:38]}) accompanied by an upward shooting hand gesture \texttt{\{\{gesture: right hand shooting diagonally upward\}\}}.\par\smallskip
\end{itemize}

\noindent\textbf{4.} \textbf{Non-verbal evaluation of financial scrutiny: When questions about unreceipted offerings, lack of frisking, or land deals are raised, how does the host/speaker react non-verbally (earnest nodding, leaning into the camera, open-palm appeals vs. eye-rolls, scoffs, dismissive hand waves, or incredulous rising pitch $\uparrow$)?}\par\smallskip
\begin{itemize}
  \item \textbf{Notes:}
Financial scrutiny is treated with \textbf{earnest, pedagogical engagement} mixed with astonished prosody at the numbers. During the chronological breakdown of the land price jump from Rs 2 crore to Rs 18.5 crore in five minutes (\texttt{[02:02]}--\texttt{[02:35]}), Sarthak leans into the desk, uses his left index finger to tap out the exact timestamps on his right hand \texttt{\{\{gesture: counting minutes on fingers\}\}}, and stresses the numerals with sharp pitch peaks: \textit{"\_panch\_ minute baad (.) \_saadhe atharah\_ crore$\uparrow$"} (\texttt{[02:27]}). When presenting the Trust's explanation (old 2011 agreement vs. prime location valuation, \texttt{[03:18]}--\texttt{[03:42]}), he shifts to a balanced two-handed weighing gesture \texttt{\{\{gesture: alternating left and right open palms like a balance scale\}\}}, visually performing fairness while still arching his eyebrows on the unresolved timing gap (\texttt{[04:39]}).\par\smallskip
\end{itemize}

\noindent\textbf{5.} \textbf{Performing certainty vs. hedging and innuendo (pauses, stops, trailing off): Are allegations delivered with definitive falling intonation $\downarrow$, finger-counting, and pointing at on-screen evidence (CCTV, documents), or hedged with hesitant pauses (.), abrupt self-interruptions (---), throat-clearing [cough], and trailing sentences (...) that imply higher-level complicity without naming it?}\par\smallskip
\begin{itemize}
  \item \textbf{Notes:}
Sarthak separates verifiable documentary facts from unproven intent using distinct prosodic contours:\par\smallskip
  \item \textbf{Documentary Facts (Certainty):} Delivered with crisp falling intonation \texttt{$\downarrow$}, finger-counting, and direct pointing at on-screen registration documents (\texttt{[02:18]}--\texttt{[02:35]}): \textit{"sham ko saat bajkar dus minute par (.) do crore mein khareedi$\downarrow$ (\texttt{[02:18]}) ... aur saat bajkar pandrah minute par (.) saadhe atharah crore ka agreement ho gaya$\downarrow$"} (\texttt{[02:27]}).
  \item \textbf{Unproven Intent / Valuation Claims (Hedging \& Skeptical Suspension):} Framed via trailing sentences \texttt{...} and open-ended rising intonation \texttt{$\uparrow$} at the close of the segment (\texttt{[04:39]}--\texttt{[05:18]}): \textit{"Ab ye waqai purane agreement ki wajah se tha (.) ya phir kuch aur..."} (\texttt{[04:45]}), accompanied by a shoulder shrug and a sideways head tilt \texttt{\{\{gesture: bilateral shoulder shrug + head tilt\}\}}. This allows him to suggest that the 5-minute flip smells suspicious without making a legally actionable accusation of corruption.
\end{itemize}

\noindent\textbf{6.} \textbf{Turn-taking, interruptions, and acoustic treatment of denials: In debates/interviews: who receives calm, uninterrupted speaking time and sympathetic nods, and who is cut off mid-sentence (---), talked over, or muted? In solo videos: how are clips of official denials framed via reaction pauses, raised eyebrows, or sarcastic chuckles?}\par\smallskip
\begin{itemize}
  \item \textbf{Notes:}
Because this is a solo YouTube explainer, Sarthak holds 100\% of the acoustic floor without panelist interruptions. However, his \textbf{multimodal framing of official denials} is highly revealing. After summarizing Champat Rai's statement that the Trust followed all legal procedures and that critics are like Gandhi's assassins (\texttt{[04:10]}--\texttt{[04:18]}), Sarthak inserts a deliberate \texttt{(1.2)} reaction pause, lets out a soft, bemused \texttt{[chuckle]}, and remarks with a wry smile: \textit{"Matlab (.) sawaal poochhne par seedha \_itna\_ bada comparison$\uparrow$"} (\texttt{[04:22]}). While he grants the Trust's substantive defence (registered bank transfers, old agreement dates, \texttt{[03:18]}--\texttt{[03:55]}) clear, uninterrupted explanation time, his post-quote facial reaction punctures the Trust's moral deflection.\par\smallskip
\end{itemize}

\noindent\textbf{7.} \textbf{Performing the boundary between religious devotion and institutional critique: How does the presenter use vocal tone, sighs, and body language to signal personal reverence for the Temple (or secular objectivity) while criticising---or defending---the Trust's financial management?}\par\smallskip
\begin{itemize}
  \item \textbf{Notes:}
Sarthak explicitly builds a verbal and non-verbal firewall between faith (\texttt{aastha}) and institutional accounting (\texttt{hisaab}) at \texttt{[04:55]}--\texttt{[05:25]}. When speaking of ordinary citizens' devotion (\textit{"crore-on logon ki aastha judi hai"\texttt{, }[04:58]\texttt{), his voice softens slightly, his pace slows, and he places his right hand lightly over his chest }\{\{gesture: right palm placed flat on chest\}\}\texttt{ with an earnest, unsmiling expression. Immediately afterward, on }"isiliye (.) \_ek-ek\_ paise ka hisaab saaf hona chahiye$\downarrow$"* (}[05:06]\texttt{), he straightens his posture and taps his index finger firmly on the desk }\{\{gesture: single firm index-finger tap on desk\}\}`. This bodily shift signals that demanding financial transparency from the Trust is an act of respect toward devotees rather than an attack on religion.\par\smallskip
\end{itemize}

\noindent\textbf{8.} \textbf{Self-presentation as the "guardian of public trust": How does the speaker use direct eye-gaze, posture, moral sighs, or emphatic gestures (pointing at the viewer, hands on heart/desk) to position themselves as the authentic protector of devotees or public accountability?}\par\smallskip
\begin{itemize}
  \item \textbf{Notes:}
Rather than posing as an angry moral crusader (like Arnab Goswami) or a grieving devotee (like Anjana Om Kashyap), Sarthak performs the role of the \textbf{calm, rational buddy-explainer} who protects the viewer from media noise. He maintains steady, relaxed eye-gaze into the lens, sits upright with open shoulders, uses inclusive pointing gestures toward the camera (\texttt{"aap bataiye comments mein"}, \texttt{[05:22]}), and lets out a light, relatable \texttt{[sigh]} of fatigue at media polarisation (\texttt{[00:34]}). His closing open-palm gesture \texttt{\{\{gesture: both palms open toward viewer\}\}} (\texttt{[05:25]}) hands the final judgment to the audience, cementing his persona as an objective, citizen-first mediator.\par\smallskip
\end{itemize}

\subsubsection*{Apply the five non-verbal features (Event 1)}

\begin{longtable}{|p{0.31\textwidth}|p{0.31\textwidth}|p{0.31\textwidth}|}
\hline
Non-verbal feature & Timestamp + spoken anchor words + non-verbal notation & What it does in context (stance, framing, implication) \\ \hline \endhead
\textbf{1. Pauses, silences \& stops} & \texttt{[02:27]} \textit{"Saat bajkar dus minute par do crore... aur (1.0) saat bajkar pandrah minute par (.) \_saadhe atharah\_ crore$\downarrow$"} & The deliberate \texttt{(1.0)} rhetorical pause before the second price creates suspense and highlights the absurdity of a 5-minute, Rs 16.5-crore valuation jump without needing shouted outrage. \\ \hline
\textbf{2. Trailing sentences (\texttt{...})} & \texttt{[04:45]} \textit{"Ab ye waqai purane agreement ki wajah se tha (.) ya phir kuch aur..."} \texttt{\{\{gesture: shoulder shrug + raised eyebrows\}\}} & Leaves the alternative explanation (\textit{"kuch aur..."}, i.e., insider profiteering) unspoken so the viewer supplies the suspicion while the host preserves his neutral explainer stance. \\ \hline
\textbf{3. Intonation, pitch \& volume} & \texttt{[02:38]} \textit{"Zameen ki keemat (.) \_ROCKET\_ ki speed se badh gayi$\uparrow$"} (marked rising pitch \texttt{$\uparrow$} + loud stress on \texttt{ROCKET}) & Uses playful prosodic hyperbole and pitch elevation to mock the economic improbability of the transaction while keeping the mood light and accessible. \\ \hline
\textbf{4. Coughs, sighs \& vocalisations} & \texttt{[04:22]} \texttt{[(1.2) pause + [chuckle]]} \textit{"Matlab (.) sawaal poochhne par seedha \_itna\_ bada comparison$\uparrow$"} & The brief ironic \texttt{[chuckle]} after quoting Champat Rai's "Gandhi's killers" remark non-verbally evaluates the Trust's rhetoric as defensive overreaction. \\ \hline
\textbf{5. Hand gestures, face \& posture} & \texttt{[03:18]}--\texttt{[03:42]} \textit{"Ek taraf aarop hai... doosri taraf Trust ka kehna hai..."} \texttt{\{\{gesture: alternating left and right open palms like a weighing scale\}\}} & Physically embodies structural balance ("both-sides" journalism), visually assuring the audience that he is weighing accusation and defence fairly. \\ \hline
\end{longtable}

\subsubsection*{Short summary (Event 1 Non-Verbal Synthesis)}

\begin{itemize}
  \item \textbf{What overall interpretation or emotional stance does the non-verbal delivery encourage?}
The non-verbal delivery encourages \textbf{skeptical, good-humoured civic scrutiny}. By pairing rapid, upbeat delivery and deadpan eyebrow raises with clear finger-counted timelines, Sarthak prevents the viewer from falling into either communal defensiveness or sensationalist panic, framing the controversy as a legitimate administrative question that deserves a transparent audit.\par\smallskip
  \item \textbf{Whose voice or physical presence is acoustically/visually foregrounded, undermined, or silenced?}
Sarthak's own mediating voice dominates the acoustic space, while both the Opposition's documentary evidence (registration timestamps) and the Trust's official rebuttals are visually foregrounded via clean on-screen graphics. However, Champat Rai's emotional counter-accusations are subtly \textbf{undermined} via reaction pauses and bemused smirks, whereas the ordinary donor's perspective is \textbf{validated} through warm, inclusive hand-to-chest gestures.\par\smallskip
  \item \textbf{What evidence supports your non-verbal reading? What complicates it (e.g., speaker baseline, format)?}
The clustering of timed pauses \texttt{(1.0)}, finger-counting, and emphatic numeral stress around the \texttt{7:10 PM} vs. \texttt{7:15 PM} registry deeds (\texttt{[02:18]}--\texttt{[02:38]}) shows where the explainer locates the core evidentiary problem. What complicates a purely ideological reading is \textbf{speaker baseline and format}: Sarthak's fast-paced jump cuts, eyebrow raises, and ironic chuckles are his trademark \textit{Sunday Show} infotainment baseline across all topics in the episode (Waqf Bill, Coaching Mafia, Floods), meaning his humour is a genre convention for youth engagement rather than targeted hostility toward the Temple Trust.\par\smallskip
\end{itemize}

\bigskip\hrule\bigskip

\subsection*{Event 2. Students' protest regarding NEET paper leak (Non-Verbal \& Paralinguistic Layer)}

\begin{itemize}
  \item \textbf{Focus:} How do prosody, pauses, trailing sentences, vocalisations, and visual body language construct the legitimacy of student protesters, police use of force, and state accountability?
  \item \textbf{Outlet:} Sarthak Goswami (\textit{Sunday Show}) | \textbf{Date:} 23 June 2024
  \item \textbf{Video title:} NEET Scam Explained! Grace Marks, Paper Leak \& Protest | \href{https://youtu.be/EM1GeZufO6E}{https://youtu.be/EM1GeZufO6E} | \textbf{Timestamp(s):} \texttt{[00:00]}--\texttt{[18:45]}
  \item \textbf{Analyst:} Baani Solanki (\texttt{2025114016}) | \textbf{Language:} Hinglish (Hindi in Roman script + English terms)
  \item \textbf{Format / Visual setup:} \texttt{[ ]} TV Panel Debate \texttt{[ ]} TV Bulletin \texttt{[x]} Solo YouTube Explainer \texttt{[ ]} Interview / Discussion
  \item \textbf{Primary speaker(s) observed:} Sarthak Goswami (solo presenter at desk using rapid jump-cuts, scorecards, statistical overlays, and brief clips of student protesters and NTA press briefings)
\end{itemize}

\subsubsection*{Event-specific non-verbal questions (Event 2)}

\noindent\textbf{1.} \textbf{Prosodic \& gestural marking of participant labels ("students/bachche" vs. "protesters" vs. "upadravi/goons"): How do vocal tone, volume, facial expression, and hand gestures shift across labels? (e.g., softened, empathetic pitch for "students" vs. harsh, staccato stress, scowls, or accusatory finger-pointing on "upadraviyon", "goons", or "anti-social elements".)}\par\smallskip
\begin{itemize}
  \item \textbf{Notes:}
Sarthak never uses criminalising labels like \textit{"upadravi"} or \textit{"goons"}. Throughout the video, he refers to the participants as \textit{"students"}, \textit{"bachche"}, or \textit{"aspirants"} (\texttt{[01:12]}, \texttt{[08:40]}, \texttt{[14:15]}) using a \textbf{softened, empathetic vocal tone}, slightly lowered pitch, and earnest furrowed brows \texttt{\{\{gesture: sympathetic brow furrow + open palms\}\}}. When describing their hard work (\textit{"do-do saal kota mein baith kar taiyari karte hain"\texttt{, }[01:28]\texttt{), his speech slows down and his voice carries genuine warmth. By contrast, when naming }"paper leak mafia"\textit{ or }"solver gang"* (}[09:45]\texttt{), his tone turns crisp and clipped, accompanied by a sharp downward chopping motion of the right hand }\{\{gesture: vertical hand chop\}\}`.\par\smallskip
\end{itemize}

\noindent\textbf{2.} \textbf{Delivery of students' stated demands: When the protesters' core demands (Education Minister's resignation, compensation, release of Sonam Wangchuk, exam reform) are mentioned, are they delivered slowly with deliberate pauses and finger-counting, rushed through in a flat monotone, mocked with air-quotes/scoffs, or trailed off (...)?}\par\smallskip
\begin{itemize}
  \item \textbf{Notes:}
The students' core demands---cancellation of fraudulent grace marks, a transparent Supreme Court-monitored probe into the Bihar/Godhra paper leaks, and complete overhaul of the NTA (\texttt{[13:50]}--\texttt{[15:10]})---are delivered \textbf{slowly and methodically with deliberate pauses and finger-counting}. Sarthak holds up three fingers one by one \texttt{\{\{gesture: sequential counting on fingers 1, 2, 3\}\}}, pausing \texttt{(0.8)} after each demand and nodding firmly with falling intonation \texttt{$\downarrow$} (*"Pehli maang (.) NTA ki accountability tay ho$\downarrow$"\texttt{, }[14:02]`), visually and acoustically validating their demands as rational and structured.\par\smallskip
\end{itemize}

\noindent\textbf{3.} \textbf{Where does vocal \& kinetic energy peak---grievance, street disruption, or electoral horse-race?: Track where the speaker's volume, pitch range, pace, and gesture intensity reach their climax: during accounts of the paper leak/student suicides, during street clashes and barricade breaches, or during partisan BJP-vs-Opposition sparring.}\par\smallskip
\begin{itemize}
  \item \textbf{Notes:}
Vocal and kinetic energy peaks unmistakably during the \textbf{statistical and administrative breakdown of student grievances and NTA anomalies} (\texttt{[03:15]}--\texttt{[07:50]}), rather than street clashes or partisan sparring:\par\smallskip
  \item \textbf{Climax 1 (Statistical Anomaly, \texttt{[03:42]}):} When revealing that \textit{"\_sadsath\_ (67) students ko full \texttt{720} out of \texttt{720} marks mile$\uparrow$!"} and that \textit{"\_ek hi\_ centre (Bahadurgarh) se \_chhah\_ toppers nikle$\uparrow$!"} (\texttt{[05:18]}), Sarthak's pitch rises sharply \texttt{$\uparrow$}, his volume surges, and he leans right into the camera lens with wide eyes and both hands spread wide in disbelief \texttt{\{\{gesture: leaning into lens + wide-eyed disbelief\}\}}.
  \item \textbf{Climax 2 (Mental Toll on Aspirants, \texttt{[15:40]}):} Energy shifts from high-volume disbelief to intense, quiet gravity \texttt{$^\circ$word$^\circ$} when discussing student depression and suicides, whereas partisan BJP-vs-Congress blame games (\texttt{[16:10]}) are brushed past rapidly with a dismissive hand wave.
\end{itemize}

\noindent\textbf{4.} \textbf{Non-verbal attribution of responsibility (sighs, stops, head shakes, pointing): When examination authorities (NTA), the Education Minister, police, opposition politicians, or students are named as responsible, what non-verbal cues accompany the attribution (exasperated sighs, mocking laughs, abrupt self-stops ---, or direct finger-pointing at the lens)?}\par\smallskip
\begin{itemize}
  \item \textbf{Notes:}
Responsibility is squarely attributed to the \textbf{National Testing Agency (NTA)} and the \textbf{Education Ministry's initial denialism} through exasperated vocalisations and mocking gestures:\par\smallskip
  \item When citing NTA's excuse of \textit{"loss of time"} and mysterious \textit{"normalization formula"} (\texttt{[06:10]}), Sarthak lets out an audible exasperated \texttt{[sigh]} followed by a mocking \texttt{[scoff]}, rolling his eyes and making physical air-quotes \texttt{\{\{gesture: air-quotes with both hands on 'loss of time'\}\}}.
  \item When referencing Education Minister Dharmendra Pradhan's initial claim that \textit{"koi paper leak ka saboot nahin hai"} (\texttt{[11:05]}) right before Bihar Police arrested the solver gang with burnt question papers (\texttt{[11:25]}), Sarthak stops mid-sentence \texttt{---}, shakes his head slowly from side to side \texttt{\{\{gesture: slow incredulous head shake\}\}}, and points at the on-screen Bihar EOU evidence report (\texttt{[11:32]}).
\end{itemize}

\noindent\textbf{5.} \textbf{Prosody of questions---open inquiry vs. loaded rhetorical challenge: Are questions about the protest delivered with level, open intonation, or with sharp rising pitch $\uparrow$, sarcastic stress, raised eyebrows, and trailing clauses (...) that presuppose protesters are violent or politically manipulated?}\par\smallskip
\begin{itemize}
  \item \textbf{Notes:}
Unlike TV studio debates that direct hostile rhetorical questions at protesters, Sarthak directs all his \textbf{loaded rhetorical challenges at the NTA and the government}. For instance, at \texttt{[04:30]}, after explaining how \texttt{718} and \texttt{719} marks are mathematically impossible under NEET's \texttt{+4 / -1} marking scheme, he asks: \textit{"NTA walon (.) ye kaunsa \_naya\_ ganit hai bhai$\uparrow$?!"} (\texttt{[04:35]}), delivered with sharp rising pitch \texttt{$\uparrow$}, sarcastic stress on \textit{"\_naya\_ ganit"}, and a classic palm-up "explain this" gesture \texttt{\{\{gesture: right arm extended, palm upturned toward screen\}\}}. Questions about the students' future (\texttt{"ab un }24\texttt{ lakh bachchon ka kya hoga$\downarrow$"}, \texttt{[17:10]}) are delivered with a falling, empathetic sigh rather than suspicion.\par\smallskip
\end{itemize}

\noindent\textbf{6.} \textbf{Visual \& paralinguistic performance of injury and violence: How are claims of violence substantiated non-verbally on screen? (Compare camera close-ups and pained vocal tone on police bandages/fractured hands vs. visual/acoustic treatment of student head injuries, lathi-charge footage, or defensive stammers/pauses when an officer or guest is asked who used force first.)}\par\smallskip
\begin{itemize}
  \item \textbf{Notes:}
Sarthak frames the crisis primarily as \textbf{institutional violence against students' futures} rather than a street riot. When showing brief B-roll of students protesting in scorching June heat and facing police barricades/detentions (\texttt{[13:20]}--\texttt{[13:45]}), his voice drops in pitch and slows down: \textit{"Chalis degree ki garmi mein (.) students sadak par utarne ko \_majboor\_ hain$\downarrow$"} (\texttt{[13:28]}). There is zero focus on police injuries or broken barricades; instead, visual close-ups highlight exhausted students holding placards and tearful parents outside exam centres (\texttt{[13:35]}).\par\smallskip
\end{itemize}

\noindent\textbf{7.} \textbf{Asymmetric delivery of police actions vs. protester actions: Compare the vocal cadence, pauses, and hand gestures used when describing police actions (e.g., calm, matter-of-fact administrative tone on "tear gas chhode gae" / "preventive arrangement") versus protester actions (alarmist pitch, striking gestures, or ironic softening as "retaliation").}\par\smallskip
\begin{itemize}
  \item \textbf{Notes:}
The asymmetry in Sarthak's video runs \textbf{in favour of the students and Bihar investigating police}, and \textbf{against central exam authorities}:\par\smallskip
  \item \textbf{Protester actions} are framed with sympathetic cadence and solidarity gestures (*"apne haq ke liye aawaz utha rahe hain"\texttt{, }[13:40]`).
  \item \textbf{Bihar Police / EOU investigation} into the Patna safe-house where burnt admit cards were recovered (\texttt{[09:50]}--\texttt{[10:40]}) is narrated with brisk, approving energy and pointing gestures at the forensic proof.
  \item \textbf{NTA's administrative actions} (releasing results 10 days early on 4 June---Election Result Day---at \texttt{[07:15]}) are narrated with a conspiratorial lean into the camera, a narrowed gaze, and a whispered aside: \textit{"Chaar June ko (.) jab poora desh election results dekh raha tha (.) $^\circ$chupke se NEET ka result nikaal diya$^\circ$"} (\texttt{[07:22]}).
\end{itemize}

\noindent\textbf{8.} \textbf{Acoustic \& visual space granted to students' own voices: Are protesting students heard in their own voice with clear audio, steady close-ups, and empathetic listening nods, or are they reduced to noisy background chaos B-roll while studio anchors/politicians speak over them?}\par\smallskip
\begin{itemize}
  \item \textbf{Notes:}
Students and educators are granted \textbf{clear, unmuted acoustic space}. At \texttt{[12:40]}--\texttt{[13:18]}, Sarthak cuts directly to clean audio bytes of NEET aspirants and educators explaining the rank inflation---where a score of \texttt{650} that previously guaranteed a government medical seat now resulted in a rank beyond \texttt{25,000}. During the transition back to studio (\texttt{[13:20]}), Sarthak nods slowly with a serious, unsmiling expression \texttt{\{\{gesture: slow, serious nod + pursed lips\}\}}, allowing the students' frustration to linger for \texttt{(1.0)} second before resuming commentary.\par\smallskip
\end{itemize}

\noindent\textbf{9.} \textbf{Non-verbal shift between the NEET paper-leak issue and the street protest: Does the speaker's non-verbal demeanour shift from sympathetic concern (sombre tone, sighs, solidarity gestures) when discussing the exam leak itself to combative, authoritarian, or dismissive delivery once students march toward Parliament?}\par\smallskip
\begin{itemize}
  \item \textbf{Notes:}
\textbf{No shift occurs.} Unlike pro-establishment TV channels that pivot from sympathy over the exam leak to hostility once students protest on the street, Sarthak maintains a continuous arc of \textbf{empathetic solidarity} from the exam hall to the street protest. He treats the street protests (\texttt{[13:20]}--\texttt{[14:30]}) as the inevitable, justified consequence of institutional deafness: \textit{"Jab sunwai nahin hogi (.) toh bachche sadak par hi aayenge na$\downarrow$"} (\texttt{[13:42]}), accompanied by an open-palm shrug of obviousness \texttt{\{\{gesture: bilateral open-palm shrug\}\}}.\par\smallskip
\end{itemize}

\subsubsection*{Apply the five non-verbal features (Event 2)}

\begin{longtable}{|p{0.31\textwidth}|p{0.31\textwidth}|p{0.31\textwidth}|}
\hline
Non-verbal feature & Timestamp + spoken anchor words + non-verbal notation & What it does in context (stance, framing, implication) \\ \hline \endhead
\textbf{1. Pauses, silences \& stops} & \texttt{[03:42]} \textit{"Is saal NEET mein (1.2) \_SADSATH\_ (\texttt{67}) students ko full \texttt{720} marks mile$\uparrow$!"} & The \texttt{(1.2)} dramatic pause before the number \texttt{67} builds anticipation and makes the statistical impossibility hit the viewer with maximum shock value. \\ \hline
\textbf{2. Trailing sentences (\texttt{...})} & \texttt{[07:28]} \textit{"Result \texttt{14} June ko aana tha, par aaya \texttt{4} June ko---jis din Lok Sabha election ke results the... ab aap khud samajh lijiye..."} & Trails off after pairing the early NEET result release with Election Day, inviting the viewer to infer deliberate media burial by the NTA without stating it as dogma. \\ \hline
\textbf{3. Intonation, pitch \& volume} & \texttt{[04:35]} \textit{"Plus char, minus ek ke hisaab se \texttt{718} aur \texttt{719} marks aa hi nahin sakte---ye kaunsa \_NAYA\_ ganit hai$\uparrow$?!"} / \texttt{[07:22]} \textit{"$^\circ$chupke se result nikaal diya$^\circ$"} & Contrasts loud, rising incredulity \texttt{$\uparrow$} on the impossible scores with a conspiratorial whispered drop \texttt{$^\circ$...$^\circ$} on the Election-Day release, dramatising both incompetence and stealth. \\ \hline
\textbf{4. Coughs, sighs \& vocalisations} & \texttt{[06:12]} \texttt{[[sigh] + [scoff]]} \textit{"Aur NTA ne iska naam diya---'grace marks for loss of time'"} & The exasperated \texttt{[sigh]} and brief \texttt{[scoff]} immediately pre-judge the NTA's "grace marks" rationale as a flimsy post-hoc cover-up before the technical rebuttal even begins. \\ \hline
\textbf{5. Hand gestures, face \& posture} & \texttt{[06:15]} \texttt{\{\{gesture: physical air-quotes on 'grace marks'\}\}} / \texttt{[07:20]} \texttt{\{\{gesture: leaning close to lens + narrowed eyes\}\}} / \texttt{[14:02]} \texttt{\{\{gesture: counting demands on 3 fingers\}\}} & Air-quotes delegitimise NTA's terminology; the camera lean-in signals an insider reveal on the \texttt{4 June} timing; finger-counting lends structured legitimacy to student demands. \\ \hline
\end{longtable}

\subsubsection*{Short summary (Event 2 Non-Verbal Synthesis)}

\begin{itemize}
  \item \textbf{What overall interpretation or emotional stance does the non-verbal delivery encourage?}
The non-verbal delivery encourages \textbf{empathetic outrage on behalf of the students combined with sharp satirical skepticism toward the NTA}. The viewer is guided to feel that \texttt{24 lakh} hardworking young citizens have been cheated by systemic incompetence and a cover-up attempt.\par\smallskip
  \item \textbf{Whose voice or physical presence is acoustically/visually foregrounded, undermined, or silenced?}
\textbf{Students, parents, and Bihar police investigators} are visually and acoustically foregrounded through clear clips, scorecards, and earnest nods. \textbf{The NTA and Education Ministry's early denials} are systematically undermined through air-quotes, scoffs, eye-rolls, and juxtaposition with burnt question-paper evidence.\par\smallskip
  \item \textbf{What evidence supports your non-verbal reading? What complicates it (e.g., speaker baseline, format)?}
The sharpest prosodic shifts---whispered delivery (\texttt{$^\circ$chupke se$^\circ$}, \texttt{[07:22]}), physical air-quotes (\texttt{[06:15]}), and \texttt{1.2}-second dramatic pauses (\texttt{[03:42]})---cluster densely around NTA's procedural anomalies. A complicating factor is that Sarthak's signature comedic meme-inserts occasionally lighten the gravity of the crisis, though he consciously suppresses jokes during the concluding section on student mental health (\texttt{[15:40]}--\texttt{[18:00]}).\par\smallskip
  \item \textbf{How do the non-verbal cues interact with the spoken words? (Do they reinforce, dramatise, soften/hedge, invert meaning via sarcasm/irony, or leave an accusation unspoken via trailing off?)}
The non-verbal cues \textbf{invert official terminology via irony} (turning \textit{"grace marks"} and \textit{"transparent process"} into objects of ridicule via air-quotes and scoffs), \textbf{dramatise statistical anomalies} via pitch peaks \texttt{$\uparrow$}, and \textbf{leave the accusation of political cover-up unspoken via trailing off} (\texttt{...}) on the \texttt{4 June} election-day result timing (\texttt{[07:28]}).\par\smallskip
  \item \textbf{What context, speaker baseline, or format factors support or complicate this interpretation? (e.g., Is this the speaker's normal style across all videos, or a marked shift on this specific actor/claim?)}
Across both Event 1 and Event 2, Sarthak's baseline is brisk, ironic, youth-centric infotainment. However, comparing Event 1 and Event 2 reveals a \textbf{marked shift in emotional investment}: whereas Event 1 (Ram Mandir land deal) stays detached and "both-sides" neutral right to the end, in Event 2 (NEET leak) his voice drops its comedic distance at \texttt{[15:40]} to adopt a sombre, protective tone toward students, reflecting his core youth/student audience demographic.\par\smallskip
  \item \textbf{How does another outlet perform the same moment, actor, or claim non-verbally?}
On \textbf{Aaj Tak} (\texttt{Event 2}), anchor Anjana Om Kashyap performs the NEET student protest with combative rising pitch \texttt{$\uparrow$}, stern finger-wagging, and repeated interruptions of student representatives while holding up close-up visuals of a police officer's bandaged hand to frame protesters as unruly (\textit{"upadravi"}). Sarthak Goswami never once wags a finger at students; his finger-pointing and sarcastic pitch peaks are aimed exclusively at the NTA's impossible \texttt{718}/\texttt{719} scorecards.\par\smallskip
  \item \textbf{My interpretation, with reasons:}
Sarthak Goswami's multimodal performance demonstrates how digital youth explainers use \textbf{paralinguistic irony (scoffs, air-quotes, whispered asides, and trailing clauses)} as a tool of institutional critique. While his verbal script stays carefully anchored in official statistics (\texttt{67} toppers, \texttt{718/719} marks, \texttt{5-minute} land registry gap), his non-verbal layer clearly signals where credibility lies---maintaining cautious neutrality when religious sensitivities are involved (Event 1), but shifting into unambiguous vocal and gestural solidarity when student futures are at stake (Event 2).\par\smallskip
\end{itemize}

\end{document}
